\subsection{Examples}

In this subsection (with the exception of Examples 7 and 8), K denotes a nondiscrete locally compact commutative field; dx denotes a Haar measure on the additive group of K.

Recall that \(\mod x = |x|\) when \(\mathbf{K} = \mathbf{R}\) , \(\mod x = |x|^2\) when \(\mathbf{K} = \mathbf{C}\) , \(\mod x = |x|_p\) when \(\mathbf{K} = \mathbf{Q}_p\) .

Example 1. --- General linear group.

Let A be the algebra \(\mathbf{M}_{n}(\mathbf{K})\) . The group \(A^{*}\) of invertible elements of A is none other than the general linear group \(\mathbf{GL}(n,\mathbf{K})\) . For every \(X\in A\) , the reduced norm \(\mathrm{Nrd}_{\mathrm{A}/\mathrm{K}}(X)\) is det X; consequently \(\mathrm{N}_{\mathrm{A}/\mathrm{K}}(X)=(\det X)^{n}\) (Alg., Ch. VIII, §12, No.~3, Prop. 8; cf.~A, III, §9, No.~3, Example 3). Since \(X\mapsto{}^{t}X\) is an isomorphism of A onto the opposite algebra,

\[
\mathrm{N} _ {\mathrm{A} ^ {0} / \mathrm{K}} (X) = \mathrm{N} _ {\mathrm{A} / \mathrm{K}} (^ {t} X) = \det (^ {t} X) ^ {n} = (\det X) ^ {n}.
\]

Then, Prop. 16 of §1, No.~11 proves that the measure

\[
\operatorname{mod} (\det X) ^ {- n} \cdot \bigotimes_ {i, j} d x _ {i j} \quad (X = (x _ {i j}))\tag{3}
\]

is a left and right Haar measure on \(\mathbf{GL}(n,\mathbf{K})\) .

To determine the relatively invariant measures on \(\mathbf{GL}(n,\mathbf{K})\) , we shall rely on the following lemma:

Lemma 5. --- The continuous representations of \(\mathbf{GL}(n,\mathbf{K})\) in \(C^{*}\) are the mappings of the form \(X \mapsto \chi(\det X)\) , where \(\chi\) is a continuous representation of \(K^{*}\) in \(C^{*}\) .

Such a mapping is obviously a continuous representation of \(\mathbf{GL}(n,\mathbf{K})\) in \(C^{*}\) . Conversely, suppose that \(\psi\) is a continuous representation of \(\mathbf{GL}(n,\mathbf{K})\) in \(C^{*}\) . For \(x \in K^{*}\) , set

\[
\widetilde {x} = \left( \begin{array}{c c c c c} x & & & & \\ & 1 & & 0 & \\ & & 1 & & \\ & 0 & & \ddots & \\ & & & & 1 \end{array} \right)
\]

and \(\chi(x) = \psi(\widetilde{x})\) . Then, for every matrix \(X \in \mathbf{GL}(n, K)\) , we have \((\det X^{-1})^{\sim} \cdot X \in \mathbf{SL}(n, K)\) . Since \(\mathbf{SL}(n, K)\) is the commutator subgroup of \(\mathbf{GL}(n, K)\) (A, III, §8, No.~9, Cor. of Prop. 17), \(\psi((\det X^{-1})^{\sim} \cdot X) = 1\) , whence

\[
\psi (X) = \psi \big ((\det X) ^ {\sim} \big) = \chi (\det X).
\]

This established, Cor. 1 of Prop. 10 of §1, No.~8 proves that the relatively invariant measures on \(\mathbf{GL}(n,\mathbf{K})\) are, up to a constant factor, the measures of the form

\[
\chi (\det X) \cdot \bigotimes_ {i j} d x _ {i j} \quad (X = (x _ {i j})),\tag{4}
\]

where \(\chi\) is a continuous representation of \(\mathbf{K}^*\) in \(\mathbf{C}^*\) .

Example 2. --- Affine group.

For every \(X \in \mathbf{GL}(n, \mathbf{K})\) and every \(x \in K^{n}\) , let \((X, x)\) be the affine linear mapping \(\xi \mapsto X\xi + x\) in \(K^{n}\) . The set of \((X, x)\) is the affine group G of \(K^{n}\) (A, II, §9, No.~4). The set T of translations is a closed normal subgroup of G, canonically isomorphic to \(K^{n}\) ; on the other hand, \(\mathbf{GL}(n, \mathbf{K})\) is a closed subgroup of G, and G is the semi-direct product of \(\mathbf{GL}(n, \mathbf{K})\) and \(T = K^{n}\) . One equips G with the (locally compact) topology for which G is the topological semi-direct product of \(\mathbf{GL}(n, \mathbf{K})\) and T (GT, III, §2, No.~10). One has

\[
(X, x) = (1, x) \cdot (X, 0).
\]

On the other hand, if \(X \in \mathbf{GL}(n, \mathbf{K})\) and \(x \in \mathbf{T}\) then, for every \(\xi \in \mathbf{K}^n\) ,

\[
(X, 0) (1, x) (X, 0) ^ {- 1} \xi = X (X ^ {- 1} \xi + x) = \xi + X x = (1, X x) \xi ,
\]

therefore the automorphism \((1,x)\mapsto(X,0)(1,x)(X,0)^{-1}\) of T has modulus mod(det X) (§1, No.~10, Prop. 15). In view of Example 1 and §2, No.~9, Remark, the measure

\[
\operatorname{mod} (\det X) ^ {- n - 1} \cdot \left(\bigotimes_ {i j} d x _ {i j}\right) \otimes \left(\bigotimes_ {i} d x _ {i}\right) \quad (X = (x _ {i j}), x = (x _ {i}))\tag{5}
\]

is a left Haar measure on G. On the other hand, by Prop. 14 of §2, No.~9,

\[
\Delta_ {\mathrm{G}} ((X, x)) = \Delta_ {\mathbf {G L} (n, K)} (X) \Delta_ {K ^ {n}} (x) (\mathrm{mod} \det X) ^ {- 1},
\]

or

\[
\Delta_ {\mathrm{G}} ((X, x)) = \operatorname{mod} (\det X ^ {- 1}).\tag{6}
\]

Thus, a right Haar measure on G is given by

\[
(\mathrm{mod} \det X) ^ {- n} \cdot \left(\bigotimes_ {i j} d x _ {i j}\right) \otimes \left(\bigotimes_ {i} d x _ {i}\right).\tag{7}
\]

Example 3. --- Strict triangular group.

Let \([1, n]\) be the set of integers \(m\) such that \(1 \leqslant m \leqslant n\) . Let \(J\) be a subset of \([1, n] \times [1, n]\) satisfying the following conditions:

\begin{enumerate}
\def\labelenumi{\arabic{enumi})}
\item
  if \((i,j)\in \mathbf{J}\) then \(i <   j\)
\item
  if \((i,j)\notin J\) then, for every integer \(k\) such that \(i < k < j\) , at least one of the two pairs \((i,k)\) and \((k,j)\) does not belong to \(J\) .
\end{enumerate}

Let \(T_J\) be the set of matrices \(Z = (z_{ij})_{1 \leqslant i \leqslant n, 1 \leqslant j \leqslant n}\) with elements in \(K\) , such that \(z_{ii} = 1\) , and \(z_{ij} = 0\) if \(i \neq j\) and \((i,j) \notin J\) . This is a closed subset of \(\mathbf{GL}(n,K)\) . The mapping \(Z \mapsto (z_{ij})_{(i,j) \in J}\) is a homeomorphism of \(T_J\) onto \(K^s\) (where \(s\) denotes the number of elements of \(J\) ). If \(Z' = (z_{ij}') \in T_J\) , then \(Z'Z = (z_{ij}'')\) with

\[
\begin{array}{c} z _ {i j} ^ {\prime \prime} = z _ {i j} + z _ {i j} ^ {\prime} + \sum_ {i <   h <   j} z _ {i h} ^ {\prime} z _ {h j} \quad \text {for} i <   j, \\ z _ {i j} ^ {\prime \prime} = 0 \text {for} i > j, \quad z _ {i i} ^ {\prime \prime} = 1, \end{array}
\]

whence \(Z^{\prime}Z \in T_{J}\) . If \(T_{J}\) is identified with \(K^{s}\) , then the mapping \(Z \mapsto Z^{\prime}Z\) (for fixed \(Z^{\prime}\) ) is identified with an affine mapping, and its determinant is 1, as one sees by ordering the pairs \((i,j) \in J\) lexicographically and applying the following lemma:

Lemma 6. --- Let L be a totally ordered finite set. For every \(\lambda \in L\) , let \(V_{\lambda}\) be a free module of finite dimension over a commutative ring \(k\) ; for \(\lambda, \mu\) in L such that \(\lambda \leqslant \mu\) , let \(f_{\lambda \mu} \in \operatorname{Hom}_k(V_\mu, V_\lambda)\) . Then the linear mapping

\[
(v _ {\lambda}) _ {\lambda \in \mathrm{L}} \mapsto \left(\sum_ {\mu \geqslant \lambda} f _ {\lambda \mu} (v _ {\mu})\right) _ {\lambda \in \mathrm{L}},
\]

from \(\prod_{\lambda \in L}V_{\lambda}\) into \(\prod_{\lambda \in L}V_{\lambda}\) , has determinant \(\prod_{\lambda \in L}\det f_{\lambda \lambda}\) .

One reduces immediately to the case that L is an interval of integers, and the lemma then follows from A, III, §8, No.~6, formula (31).

If \(Z \in \mathrm{T}_{\mathrm{J}}\) , one then sees that there exists \(Z' \in \mathrm{T}_{\mathrm{J}}\) such that \(Z'Z = I_n\) , whence \(Z' = Z^{-1}\) . Thus, \(\mathrm{T}_{\mathrm{J}}\) is a closed subgroup of \(\mathbf{GL}(n,\mathrm{K})\) . On the other hand, Prop. 15 of §1, No.~10 shows that the measure

\[
\bigotimes_ {(i, j) \in J} d z _ {i j}
\]

is a left Haar measure on \(T_{J}\) . By calculating \(ZZ'\) one sees in the same way that this measure is a right Haar measure on \(T_{J}\) .

There is an analogous result if, in the definition of \(T_{J}\) , the roles of rows and columns are interchanged.

When J is the set of pairs \((i,j)\) such that i\textless j, the group \(T_{J}\) is called the upper strict triangular group of order n over K, and is denoted \(\mathrm{T}_{1}(n,\mathrm{K})\) . Its transpose is called the lower strict triangular group.

Example 4. --- Large triangular group.

Let \(n_1, \ldots, n_r\) be integers \(\geqslant 1\) . Set \(p_k = n_1 + \ldots + n_{k-1}\) and \(n = p_{r+1} = n_1 + \cdots + n_r\) . Let \(I_k\) be the set of integers \(j\) such that \(p_k < j \leqslant p_{k+1}\) , and \(J\) the union of the \(I_k \times I_l\) for \(k < l\) . Let \(G\) be the closed subgroup of \(\mathbf{GL}(n, K)\) whose elements are the matrices \((Z_{kl})_{1 \leqslant k \leqslant r, 1 \leqslant l \leqslant r}\) such that:

\begin{enumerate}
\def\labelenumi{\arabic{enumi})}
\item
  each \(Z_{kl}\) is a matrix \((z_{ij})_{i\in \mathrm{I}_k,j\in \mathrm{I}_l}\) of elements of \(\mathbf{K}\) , with \(n_k\) rows and \(n_l\) columns;
\item
  \(Z_{kl} = 0\) for \(k > l\) ;
\item
  \(Z_{kk} \in \mathbf{GL}(n_k, \mathrm{K})\) for \(1 \leqslant k \leqslant r\) .
\end{enumerate}

The formula for block multiplication

\[
\begin{array}{r} \left( \begin{array}{c c c c} Z _ {1 1} & 0 & \ldots & 0 \\ 0 & Z _ {2 2} & \ldots & 0 \\ \vdots & \vdots & \ddots & \vdots \\ 0 & 0 & \ldots & Z _ {r r} \end{array} \right) \left( \begin{array}{c c c c} 1 & Z _ {1 2} & \ldots & Z _ {1 r} \\ 0 & 1 & \ldots & Z _ {2 r} \\ \vdots & \vdots & \ddots & \vdots \\ 0 & 0 & \ldots & 1 \end{array} \right) \\ = \left( \begin{array}{c c c c} Z _ {1 1} & Z _ {1 1} Z _ {1 2} & \ldots & Z _ {1 1} Z _ {1 r} \\ 0 & Z _ {2 2} & \ldots & Z _ {2 2} Z _ {2 r} \\ \vdots & \vdots & \ddots & \vdots \\ 0 & 0 & \ldots & Z _ {r r} \end{array} \right) \end{array}\tag{8}
\]

shows that G is the topological semi-direct product of the subgroup D of elements \((Z_{kl}) \in G\) such that \(Z_{kl} = 0\) for \(k \neq l\) and the subgroup \(T_{J}\) of Example 3. Moreover, D is isomorphic to the direct product of the groups \(\mathbf{GL}(n_{k}, \mathbf{K})\) for \(1 \leqslant k \leqslant r\) .

Let \(J'\) be the set of pairs \((j, i)\) for \((i, j) \in J\) and let H be the set of pairs \((i, j) \in [1, n] \times [1, n]\) not belonging to \(J'\) . Let \(Z' = (z_{ij})_{1 \leqslant i \leqslant n, 1 \leqslant j \leqslant n}\) be an element of G. By Prop. 14 of §2, No.~9 and the above Examples 1 and 3, one obtains a left Haar measure on G by taking the image of the measure

\[
\bigotimes_ {k = 1} ^ {r} \left(\left(\operatorname{mod} \det Z _ {k k}\right) ^ {- n _ {k}} \cdot \bigotimes_ {i, j \in I _ {k}} d z _ {i j}\right) \otimes \left(\bigotimes_ {(i, j) \in J} d z _ {i j}\right)
\]

under the mapping

\[
\bigl ((Z _ {k k}), (Z _ {k l}) \bigr) \mapsto \left( \begin{array}{c c c c} Z _ {1 1} & Z _ {1 1} Z _ {1 2} & \ldots & Z _ {1 1} Z _ {1 r} \\ 0 & Z _ {2 2} & \ldots & Z _ {2 2} Z _ {2 r} \\ \vdots & \vdots & \ddots & \vdots \\ 0 & 0 & \ldots & Z _ {r r} \end{array} \right).
\]

Now, consider, for \(k < l\) , the vector space of matrices \(Z_{kl} = (z_{ij})_{i \in \mathbf{I}_k, j \in \mathbf{I}_l}\) . It is the direct sum of the \(n_l\) subspaces \(\mathbf{M}_j\) ( \(j \in \mathbf{I}_l\) ) formed by the matrices such that \(z_{ih} = 0\) for \(h \neq j\) . Each of these subspaces \(\mathbf{M}_j\) is stable under the mapping \(Z_{kl} \mapsto Z_{kk}Z_{kl}\) , and the restriction of this mapping to \(\mathbf{M}_j\) has matrix \(Z_{kk}\) . Consequently (§1, No.~10, Prop. 15) the image of the measure

\(\bigotimes_{i \in \mathbf{I}_k, j \in \mathbf{I}_l} dz_{ij}\) under the mapping \(Z_{kl} \mapsto Z_{kk} Z_{kl}\) is

\[
(\mathrm{mod} \det Z _ {k k}) ^ {- n _ {l}} \cdot \bigotimes_ {i \in \mathrm{I} _ {k}, j \in \mathrm{I} _ {l}} d z _ {i j}.
\]

A left Haar measure on G is therefore given by

\[
\prod_ {k = 1} ^ {r} (\mathrm{mod} \det Z _ {k k}) ^ {- q _ {k}} \cdot \bigotimes_ {(i, j) \in \mathrm{H}} d z _ {i j}\tag{9}
\]

with \(q_{k} = \sum_{k\leqslant l\leqslant r}n_{l} = n - p_{k}\)

Let us calculate the modulus of G, again using Prop. 14 of §2. The groups D and T \(_{J}\) are unimodular; on the other hand:

\[
\begin{array}{c} \left( \begin{array}{c c c c} Z _ {1 1} & 0 & \ldots & 0 \\ 0 & Z _ {2 2} & \ldots & 0 \\ \vdots & \vdots & \ddots & \vdots \\ 0 & 0 & \ldots & Z _ {r r} \end{array} \right) \left( \begin{array}{c c c c} 1 & Z _ {1 2} & \ldots & Z _ {1 r} \\ 0 & 1 & \ldots & Z _ {2 r} \\ \vdots & \vdots & \ddots & \vdots \\ 0 & 0 & \ldots & 1 \end{array} \right) \left( \begin{array}{c c c c} Z _ {1 1} & 0 & \ldots & 0 \\ 0 & Z _ {2 2} & \ldots & 0 \\ \vdots & \vdots & \ddots & \vdots \\ 0 & 0 & \ldots & Z _ {r r} \end{array} \right) ^ {- 1} \\ = \left( \begin{array}{c c c c} 1 & Z _ {1 2} ^ {\prime} & \ldots & Z _ {1 r} ^ {\prime} \\ 0 & 1 & \ldots & Z _ {2 r} ^ {\prime} \\ \vdots & \vdots & \ddots & \vdots \\ 0 & 0 & \ldots & 1 \end{array} \right), \end{array}
\]

where \(Z_{kl}^{\prime}=Z_{kk}Z_{kl}Z_{ll}^{-1}\) . Taking into account Example 3, and Prop. 15 of §1, No.~10, and arguing as above, one sees that if \(X=\mathrm{diag}(Z_{11},\ldots,Z_{rr})\in\mathbf{D}\) then the modulus of the automorphism \(Z\mapsto X^{-1}ZX\) of \(T_{J}\) is

\[
\prod_ {k <   l} (\mathrm{mod} \det Z _ {k k}) ^ {- n _ {l}} (\mathrm{mod} \det Z _ {l l}) ^ {n _ {k}},
\]

therefore

\[
\Delta_ {\mathrm{G}} (Z) = \prod_ {k = 1} ^ {r} (\mathrm{mod} \det Z _ {k k}) ^ {n + n _ {k} - 2 q _ {k}}.\tag{10}
\]

The transposed group \(G'\) of G is studied in the same manner. For \(G'\) , one finds as left Haar measure

\[
\prod_ {k = 1} ^ {r} (\mathrm{mod} \det Z _ {k k}) ^ {- p _ {k + 1}} \cdot \bigotimes_ {(j, i) \in \mathrm{H}} d z _ {i j},
\]

and as modulus

\[
\prod_ {k = 1} ^ {r} (\mathrm{mod} \det Z _ {k k}) ^ {n + n _ {k} - 2 p _ {k + 1}}.
\]

If in particular one takes \(n_{1}=\ldots=n_{r}=1\) , one finds as group G the group \(T(n,K)^{*}\) of invertible elements of the subalgebra of \(\mathbf{M}_{n}(\mathbf{K})\) formed by the matrices \(X=(x_{ij})\) such that \(x_{ij}=0\) for i\textgreater j. This algebra, which we shall denote \(\mathrm{T}(n,\mathrm{K})\) , is called the upper triangular algebra, and the group \(\mathrm{T}(n,\mathrm{K})^{*}\) is called the upper large triangular group of order n over K. The preceding formulas then take the following form: a left Haar measure on \(\mathrm{T}(n,\mathrm{K})^{*}\) is

\[
\prod_ {i = 1} ^ {n} (\mathrm{mod} z _ {i i}) ^ {i - n - 1} \cdot \bigotimes_ {i \leqslant j} d z _ {i j} \quad (Z = (z _ {i j}))\tag{9 bis}
\]

and the modulus of \(\mathrm{T}(n,\mathbf{K})^*\) is

(10 bis)

\[
\Delta_ {\mathrm{T} (n, \mathrm{K}) ^ {*}} (Z) = \prod_ {i = 1} ^ {n} (\mathrm{mod} z _ {i i}) ^ {2 i - n - 1} \quad \left(Z = \left(z _ {i j}\right)\right).
\]

For the transpose of \(\mathrm{T}(n,\mathbf{K})^*\) , or lower large triangular group, one finds as left Haar measure

\[
\prod_ {i = 1} ^ {n} (\mathrm{mod} z _ {i i}) ^ {- i} \cdot \bigotimes_ {i \geqslant j} d z _ {i j},
\]

and as modulus

\[
\prod_ {i = 1} ^ {n} (\mathrm{mod} z _ {i i}) ^ {n + 1 - 2 i}.
\]

Remark. --- The group \(\mathrm{T}(n,\mathrm{K})^{*}\) is a closed subgroup of \(\mathbf{GL}(n,\mathrm{K})\) , and \(\Delta_{\mathrm{T}(n,\mathrm{K})^{*}}\left((z_{ij})\right)=\prod_{i=1}^{n}(\mathrm{mod}z_{ii})^{2i-n-1}\) . We saw in Example 1 that \(\Delta_{\mathbf{GL}(n,\mathrm{K})}=1\) . If n\textgreater1, the function

\[
\Delta_ {\mathrm{T} (n, \mathrm{K}) ^ {*}} / \Delta_ {\mathbf {G L} (n, \mathrm{K})}
\]

on \(\mathrm{T}(n,\mathbf{K})^*\) cannot be extended to a continuous representation of \(\mathbf{GL}(n,\mathbf{K})\) in \(\mathbf{C}^*\) (because such a representation would be equal to 1 on \(\mathbf{SL}(n,\mathbf{K})\) by Lemma 5, whereas \(\mathrm{mod}(z_{11})^{1 - n}\neq 1\) for \(z_{11}\) suitably chosen). It follows that the homogeneous space \(\mathbf{GL}(n,\mathbf{K}) / \mathrm{T}(n,\mathbf{K})^*\) admits no relatively invariant measure if \(n > 1\) ( \(\S 2\) , No.~6, Cor. 1 of Th. 3).

This homogeneous space may be identified, for n = 2, with the projective line over K. For, let \((e_{1}, e_{2})\) be the canonical basis of \(K^{2}\) . The group \(\mathbf{GL}(2, \mathbf{K})\) operates transitively on the set of lines of \(K^{2}\) with 0 omitted, and the stabilizer of \(Ke_{1} - \{0\}\) is \(T(2, K)^{*}\) .

Example 5. --- Special triangular group.

Let us take up again the notations at the beginning of Example 4, and consider the subgroup \(\mathbf{G}_{1} = \mathbf{G} \cap \mathbf{SL}(n, \mathbf{K})\) . This subgroup is the topological semi-direct product of the group \(\mathbf{D}_{1} = \mathbf{D} \cap \mathbf{SL}(n, \mathbf{K})\) with \(T_{J}\) . The group \(D_{1}\) has a normal subgroup A isomorphic to \(\mathbf{SL}(n_{r}, \mathbf{K})\) , namely the subgroup consisting of the elements \(\text{diag}(Z_{kk})\) with \(Z_{kk} = 1\) for k \textless{} r. The homomorphism

\[
\varphi : \operatorname{diag} (Z _ {1 1}, \dots , Z _ {r r}) \mapsto (Z _ {1 1}, \dots , Z _ {r - 1, r - 1})
\]

of \(D_{1}\) into \(\mathbf{GL}(n_{1},\mathbf{K})\times\cdots\times\mathbf{GL}(n_{r-1},\mathbf{K})\) is surjective and has kernel A. On the other hand, \(\varphi\) is continuous. Taking into account Lemma 2 of Appendix I, \(D_{1}/A\) may be identified with \(\mathbf{GL}(n_{1},\mathbf{K})\times\cdots\times\mathbf{GL}(n_{r-1},\mathbf{K})\) . We shall denote by \(\mu\) the Haar measure of A (cf.~Example 6) and by

\[
\alpha = \bigotimes_ {k = 1} ^ {r - 1} \left(\left(\operatorname{mod} \det Z _ {k k}\right) ^ {- n _ {k}} \cdot \bigotimes_ {i, j \in I _ {k}} d z _ {i j}\right) \otimes^ {\prime} d \mu \left(Z _ {r r}\right)
\]

the Haar measure on \(\mathbf{D}_1\) such that

\[
\alpha / \mu = \bigotimes_ {k = 1} ^ {r - 1} \left(\left(\operatorname{mod} \det Z _ {k k}\right) ^ {- n _ {k}} \cdot \bigotimes_ {i, j \in I _ {k}} d z _ {i j}\right)
\]

(§2, No.~7, Prop. 10). One then shows as in Example 4 that a left Haar measure on \(\mathbf{G}_1\) is given by

\[
\begin{array}{l} \bmod \Big (\prod_ {k = 1} ^ {r - 1} (\det Z _ {k k}) ^ {n _ {k} - q _ {k}} \Big) \\ \qquad \cdot \left[ \bigotimes_ {k = 1} ^ {r - 1} \big ((\bmod \det Z _ {k k}) ^ {- n _ {k}} \cdot \bigotimes_ {i, j \in I _ {k}} d z _ {i j} \big) \otimes^ {\prime} d \mu (Z _ {r r}) \right] \otimes \bigotimes_ {(i, j) \in J} d z _ {i j}. \end{array}
\]

Since \(\mathbf{G}_1\) is normal in \(\mathbf{G}\) , the modulus of \(\mathbf{G}_1\) is the restriction of that of \(\mathbf{G}\) (§2, No.~7, Prop. 10 b)).

If \(n_{r}=1\) , the subgroup A reduces to the neutral element, and a left Haar measure on G is

\[
\operatorname{mod} \left(\prod_ {k = 1} ^ {r - 1} (\det Z _ {k k}) ^ {- q _ {k}}\right) \cdot \bigotimes_ {k = 1} ^ {r - 1} \left(\bigotimes_ {i, j \in I _ {k}} d z _ {i j}\right) \otimes \bigotimes_ {(i, j) \in J} d z _ {i j}.
\]

If one takes \(n_{1}=n_{2}=\ldots=n_{r}=1\) , the group \(G_{1}\) obtained is called the upper special triangular group and its transpose \(G_{1}^{\prime}\) is called the lower special triangular group. A left Haar measure on \(G_{1}\) is

\[
\operatorname{mod} \left(\prod_ {i = 1} ^ {n - 1} z _ {i i} ^ {i - n - 1}\right) \cdot \left(\bigotimes_ {i = 1} ^ {n - 1} d z _ {i i}\right) \otimes \left(\bigotimes_ {1 \leqslant i <   j \leqslant n} d z _ {i j}\right)\tag{11}
\]

and the modulus of \(\mathbf{G}_1\) is

\[
\operatorname{mod} \left(\prod_ {i = 1} ^ {n - 1} z _ {i i} ^ {2 i - 2 n}\right).\tag{12}
\]

For \(G_{1}^{\prime}\) one finds similarly the left Haar measure

\[
\operatorname{mod} \left(\prod_ {i = 1} ^ {n - 1} z _ {i i} ^ {n - i - 1}\right) \cdot \left(\bigotimes_ {i = 1} ^ {n - 1} d z _ {i i}\right) \otimes \left(\bigotimes_ {1 \leqslant j <   i \leqslant n} d z _ {i j}\right)
\]

and modulus

\[
\operatorname{mod} \biggl (\prod_ {i = 1} ^ {n - 1} z _ {i i} ^ {2 n - 2 i} \biggr).
\]

Example 6. --- Special linear group.

The closed subgroups \(\mathrm{T}_{1}(n,\mathrm{K})\) and \({}^{t}(\mathrm{T}(n,\mathrm{K})^{*})\) of \(\mathbf{GL}(n,\mathbf{K})\) have intersection \(\{e\}\) . Thus the mapping \((M,N)\mapsto M\cdot N\) is a continuous bijection \(\varphi\) of \(\mathrm{T}_{1}(n,\mathrm{K})\times{}^{t}(\mathrm{T}(n,\mathrm{K})^{*})\) onto a subset \(\Omega\) of \(\mathbf{GL}(n,\mathbf{K})\) .

Lemma 7. --- a) Let \(U = (u_{ij}) \in \mathbf{GL}(n, K)\) . In order that \(U \in \Omega\) , it is necessary and sufficient that \(\det(u_{ij})_{k \leqslant i,j \leqslant n} \neq 0\) for \(k = 2, 3, \ldots, n\) .

\begin{enumerate}
\def\labelenumi{\alph{enumi})}
\setcounter{enumi}{1}
\item
  \(\Omega\) is an open subset of \(\mathbf{GL}(n,\mathbf{K})\) .
\item
  The mapping \(\varphi\) is a homeomorphism of \(\mathrm{T}_1(n,\mathrm{K})\times {}^t (\mathrm{T}(n,\mathrm{K})^*)\) onto \(\Omega\) .
\end{enumerate}

In order that \(U \in \Omega\) , it is necessary and sufficient that there exist a \(Z = (z_{ij}) \in \mathrm{T}_1(n,\mathrm{K})\) such that \(ZU \in {}^t (\mathbf{T}(n,\mathbf{K}))\) (then necessarily

\(ZU \in {}^{t}(\mathrm{T}(n,\mathrm{K})^{*})\) since \(U\) and \(Z\) are invertible). By what we saw earlier, if \(Z\) exists then it is unique. Thus, in order that \(U \in \Omega\) , it is necessary and sufficient that the linear system

\[
\sum_ {k = 1} ^ {n} z _ {i k} u _ {k j} = 0 \quad (1 \leqslant i <   j \leqslant n)
\]

(where \((z_{ij})\in \mathrm{T}_1(n,\mathbf{K})\) ) admit a unique solution. Now, this system may be written

\[
\sum_ {k = i + 1} ^ {n} z _ {i k} u _ {k j} = - u _ {i j} \quad (1 \leqslant i <   j \leqslant n).\tag{13}
\]

For fixed \(i\) , one has a system of \(n - i\) equations in the unknowns \(z_{i,i+1}\) , \(z_{i,i+2}\) , \ldots, \(z_{i,n}\) ; for these systems to admit unique solutions, it is necessary and sufficient that

\[
\det (u _ {k j}) _ {i + 1 \leqslant k \leqslant n, i + 1 \leqslant j \leqslant n} \neq 0
\]

for \(i = 1,2,\ldots ,n - 1\) . This proves a). From this it follows that \(\Omega\) is open in \(\mathbf{GL}(n,\mathbf{K})\) . On the other hand, on solving the system (13) by means of Cramer's formulas, the \(z_{ij}\) are obtained as rational functions of the \(u_{ij}\) with nonzero denominators in \(\Omega\) , therefore \(Z\) depends continuously on \(U\) in \(\Omega\) , which proves c).

Now let \(\mathbf{G}_{1}^{\prime}\subset{}^{t}(\mathrm{T}(n,\mathrm{K})^{*})\) be the lower special triangular group. The mapping \((M,N)\mapsto M\cdot N\) is a continuous bijection \(\psi\) of \(\mathrm{T}_{1}(n,\mathrm{K})\times\mathbf{G}_{1}^{\prime}\) onto a subset \(\Omega^{\prime}\) of \(\mathbf{SL}(n,\mathbf{K})\) .

Lemma 8. --- a) Let \(U = (u_{ij}) \in \mathbf{SL}(n, \mathbf{K})\) . In order that \(U \in \Omega'\) , it is necessary and sufficient that \(\det(u_{ij})_{k \leqslant i,j \leqslant n} \neq 0\) for \(k = 2,3,\ldots,n\) .

\begin{enumerate}
\def\labelenumi{\alph{enumi})}
\setcounter{enumi}{1}
\item
  \(\Omega'\) is an open subset of \(\mathbf{SL}(n,\mathbf{K})\) .
\item
  The mapping \(\psi\) is a homeomorphism of \(\mathrm{T}_1(n,\mathbf{K})\times \mathrm{G}_1'\) onto \(\Omega^{\prime}\) . For, let \(M\in \mathrm{T}_1(n,\mathbf{K})\) and \(N\in {}^t (\mathrm{T}(n,\mathbf{K})^*)\) . In order that \(M\cdot N\in \mathbf{SL}(n,\mathbf{K})\) , it is necessary and sufficient that \(N\in \mathbf{G}_1'\) . Therefore \(\Omega^{\prime} = \mathbf{SL}(n,\mathbf{K})\cap \Omega\) and Lemma 8 follows at once from Lemma 7.
\end{enumerate}

PROPOSITION 6. --- a) The group \(\mathbf{SL}(n,\mathbf{K})\) is unimodular.

\begin{enumerate}
\def\labelenumi{\alph{enumi})}
\setcounter{enumi}{1}
\item
  Let \(\mu_1\) and \(\mu_2\) be left Haar measures on the upper strict triangular group \(\mathrm{T}_1(n, \mathrm{K})\) and the lower special triangular group \(\mathrm{G}_1'\) , respectively. The image of \(\mu_1 \otimes \mu_2\) under the homeomorphism \((M, N) \mapsto M \cdot N^{-1}\) of \(\mathrm{T}_1(n, \mathrm{K}) \times \mathrm{G}_1'\) onto \(\Omega'\) is the restriction to \(\Omega'\) of a Haar measure on \(\mathrm{SL}(n, \mathrm{K})\) .
\item
  The complement of \(\Omega'\) in \(\mathbf{SL}(n,\mathbf{K})\) is negligible for the Haar measure of \(\mathbf{SL}(n,\mathbf{K})\) .
\end{enumerate}

The group \(\mathbf{GL}(n,\mathbf{K})\) is unimodular (Example 1), and \(\mathbf{SL}(n,\mathbf{K})\) is a normal subgroup of \(\mathbf{GL}(n,\mathbf{K})\) , hence is unimodular ( \(\S 2\) , No.~7, Prop. 10 b)). The assertion b) follows from a), Lemma 8, and Prop. 13 of \(\S 2\) , No.~9. Let us prove c). By Lemma 8 a), it suffices to prove the following: if \(p((u_{ij})_{1 \leqslant i,j \leqslant n})\) is a polynomial, not identically zero on \(\mathbf{SL}(n,\mathbf{K})\) , then the set E of \(U \in \mathbf{SL}(n,\mathbf{K})\) such that \(p(U) = 0\) is negligible for the Haar measure. Taking into account \(\S 1\) , No.~10, Cor. of Prop. 13, the topology of \(\mathbf{SL}(n,\mathbf{K})\) has a countable base. It therefore suffices to prove that for every \(U_0 \in \mathbf{E}\) , there exists a neighborhood of \(U_0\) in \(\mathbf{SL}(n,\mathbf{K})\) whose intersection with E is negligible; or again that there exists a neighborhood W of I in \(\mathbf{SL}(n,\mathbf{K})\) such that \(U_0^{-1}\mathbf{E} \cap \mathbf{W}\) is negligible. Let us take \(\mathbf{W} = \Omega'\) . In view of b), it all comes down to showing that the set of pairs \((M,N) \in \mathrm{T}_1(n,\mathbf{K}) \times \mathrm{G}_1'\) such that \(p(U_0MN) = 0\) is negligible for \(\mu_1 \otimes \mu_2\) . By the expressions for \(\mu_1\) and \(\mu_2\) (calculated in Examples 3 and 5), this will result from the following lemma:

Lemma 9. --- Let \(\psi\) be a polynomial \(\neq 0\) of \(\mathbf{K}[\mathbf{X}_1, \ldots, \mathbf{X}_r]\) . In the space \(\mathbf{K}^r\) , the set \(\mathbf{N}\) defined by \(\psi(x_1, \ldots, x_r) = 0\) is negligible for the Haar measure.

Let us argue by induction on r. The lemma is evident for r = 1, since N is then a finite set. Changing if necessary the numbering of the variables, we can suppose that \(\psi \notin K[X_{1}, \ldots, X_{r-1}]\) ; write

\[
\psi \left(\mathrm{X} _ {1}, \dots , \mathrm{X} _ {r}\right) = \mathrm{X} _ {r} ^ {m} \psi_ {0} \left(\mathrm{X} _ {1}, \dots , \mathrm{X} _ {r - 1}\right) + \dots + \psi_ {m} \left(\mathrm{X} _ {1}, \dots , \mathrm{X} _ {r - 1}\right)
\]

with \(m > 0\) and \(\psi_0 \neq 0\) . In the space \(\mathbf{K}^{r-1}\) , let \(\mathbf{N}_0\) be the set defined by \(\psi_0(x_1, \ldots, x_{r-1}) = 0\) , which is negligible by the induction hypothesis. For every \((x_1, \ldots, x_{r-1}) \notin \mathbf{N}_0\) , the set of \(x_r \in \mathbf{K}\) such that \((x_1, \ldots, x_{r-1}, x_r) \in \mathbf{N}\) is finite, therefore negligible. Since \(\mathbf{K}^r\) is countable at infinity (§1, No.~10, Cor. of Prop. 13), \(\mathbf{N} \cap [(\mathbf{K}^{r-1} - \mathbf{N}_0) \times \mathbf{K}]\) is negligible in \(\mathbf{K}^r\) (Ch. V, §8, No.~2, Prop. 4). Therefore \(\mathbf{N}\) is negligible.

Example 7. --- Iwasawa decomposition of \(\mathbf{GL}(n,\mathbf{K})\).

In this example, K denotes one of the fields R, C, H. If \(\lambda \in K\) , \(\overline{\lambda}\) is defined to be equal to \(\lambda\) if K = R, and to the conjugate of \(\lambda\) if K = C or H. Let E be a right vector space over K of dimension n, and let \(\Phi\) be a nondegenerate positive hermitian form on E.

Lemma 10. --- Let \((f_1, f_2, \ldots, f_n)\) be a basis of E.

\begin{enumerate}
\def\labelenumi{\alph{enumi})}
\item
  There exists one and only one orthonormal basis \((e_{1}, e_{2}, \ldots, e_{n})\) of E such that \(f_{i} = e_{1}\alpha_{i1} + e_{2}\alpha_{i2} + \cdots + e_{i}\alpha_{ii}\) ( \(i = 1, 2, \ldots, n\) ) with \(\alpha_{ii} > 0\) for all i.
\item
  For fixed \(\Phi\) , the \(e_i\) and \(\alpha_{ij}\) depend continuously on \((f_1, \ldots, f_n) \in \mathbf{E}^n\) .
\end{enumerate}

Let \(E_{i}=f_{1}K+f_{2}K+\cdots+f_{i}K\) , which has dimension i. Let \(g_{i}\) be a nonzero element of \(E_{i}\) orthogonal to \(E_{i-1}\) and such that \(\Phi(g_{i},g_{i})=1\) . By induction on i, one sees that \((g_{1},\ldots,g_{i})\) is an orthonormal basis of \(E_{i}\) . In particular, \((g_{1},\ldots,g_{n})\) is an orthonormal basis of E. Let \(\lambda_{i}=\Phi(f_{i},g_{i})\) . Since \(f_{i}\notin E_{i-1}\) , one has \(\lambda_{i}\neq0\) . Set \(e_{i}=g_{i}\left|\lambda_{i}\right|\lambda_{i}^{-1}\) . Then

\[
\Phi (e _ {i}, e _ {i}) = | \lambda_ {i} | ^ {2} \overline {{{{\lambda}}}} _ {i} ^ {- 1} \Phi (g _ {i}, g _ {i}) \lambda_ {i} ^ {- 1} = 1,
\]

thus \((e_1, \ldots, e_i)\) is also an orthonormal basis of \(\mathbf{E}_i\) ; moreover, \(\Phi(e_i, f_i) = |\lambda_i| \overline{\lambda}_i^{-1} \Phi(g_i, f_i) = |\lambda_i| > 0\) , thus the \(e_i\) have the properties of a). Let \((e_1', \ldots, e_n')\) be another orthonormal basis of \(\mathbf{E}\) with the same properties. One sees by induction on \(i\) that \((e_1', \ldots, e_i')\) must be a basis of \(\mathbf{E}_i\) , therefore \(e_i' = e_i \mu_i\) for some \(\mu_i \in \mathbf{K}\) . Then

\[
1 = \Phi (e _ {i} ^ {\prime}, e _ {i} ^ {\prime}) = \overline {{\mu}} _ {i} \Phi (e _ {i}, e _ {i}) \mu_ {i} = \overline {{\mu}} _ {i} \mu_ {i},
\]

and \(0 < \Phi(e_i', f_i) = \overline{\mu}_i \Phi(e_i, f_i)\) , therefore \(\mu_i > 0\) and \(\mu_i^2 = 1\) , thus \(\mu_i = 1\) , whence a). Suppose already proven that the \(e_i\) and \(\alpha_{ij}\) depend continuously on \((f_1, \ldots, f_n)\) for \(i < i_0\) , and let us prove that \(e_{i_0}\) and the \(\alpha_{i_0j}\) depend continuously on \((f_1, \ldots, f_n)\) . For \(j < i_0\) , \(\overline{\alpha}_{i_0j} = \Phi(f_{i_0}, e_j)\) depends continuously on \((f_1, \ldots, f_n)\) by the induction hypothesis. On the other hand,

\[
\Phi (f _ {i _ {0}}, f _ {i _ {0}}) = | \alpha_ {i _ {0} 1} | ^ {2} + | \alpha_ {i _ {0} 2} | ^ {2} + \dots + | \alpha_ {i _ {0}, i _ {0} - 1} | ^ {2} + \alpha_ {i _ {0} i _ {0}} ^ {2},
\]

thus \(\alpha_{i_0i_0}\) depends continuously on \((f_1,\ldots ,f_n)\) . Therefore

\[
e _ {i _ {0}} = (f _ {i _ {0}} - e _ {1} \alpha_ {i _ {0} 1} - \dots - e _ {i _ {0} - 1} \alpha_ {i _ {0}, i _ {0} - 1}) \alpha_ {i _ {0} i _ {0}} ^ {- 1}
\]

depends continuously on \((f_1, \ldots, f_n)\) .

Henceforth let \(\mathbf{E} = \mathbf{K}^n\) and let us take for \(\Phi\) the form

\[
\overline {{{x}}} _ {1} y _ {1} + \dots + \overline {{{x}}} _ {n} y _ {n}.
\]

Recall that \(\mathbf{U}(n,\mathbf{K})\) denotes the corresponding unitary group. Even when K is noncommutative, we shall again denote by \(\mathrm{T}_{1}(n,\mathrm{K})\) the group of upper triangular matrices of \(\mathbf{M}_{n}(\mathrm{K})\) whose diagonal entries are all 1.

PROPOSITION 7. --- Let \(D_{+}^{*}\) be the group of diagonal matrices with diagonal elements \textgreater0. The mapping \((U,D,T)\mapsto UDT\) is a homeomorphism of \(\mathbf{U}(n,\mathbf{K})\times\mathbf{D}_{+}^{*}\times\mathbf{T}_{1}(n,\mathbf{K})\) onto \(\mathbf{GL}(n,\mathbf{K})\) .

Let \((\varepsilon_1, \ldots, \varepsilon_n)\) be the canonical basis of \(\mathbf{K}^n\) . Let \(X \in \mathbf{GL}(n, \mathbf{K})\) . Then the \(X \cdot \varepsilon_i = f_i\) form a basis of \(\mathbf{E}\) . To this basis \((f_i)\) one can associate a basis \((e_i)\) as in Lemma 10. Let \(U\) be the matrix of the unitary automorphism of \(\mathbf{E}\) that transforms \(\varepsilon_i\) into \(e_i\) . Then

\[
U ^ {- 1} \cdot f _ {i} = \varepsilon_ {1} \alpha_ {i 1} + \varepsilon_ {2} \alpha_ {i 2} + \dots + \varepsilon_ {i} \alpha_ {i i}
\]

with \(\alpha_{ii} > 0\) for \(i = 1,2,\ldots ,n\) . Thus \(X = UC\) , where \(C\) is the matrix

\[
\left( \begin{array}{c c c c} \alpha_ {1 1} & \alpha_ {2 1} & \ldots & \alpha_ {n 1} \\ 0 & \alpha_ {2 2} & \ldots & \alpha_ {n 2} \\ \vdots & \vdots & \ddots & \vdots \\ 0 & 0 & \ldots & \alpha_ {n n} \end{array} \right).
\]

Moreover, U and C depend continuously on X by Lemma 10. On the other hand, formula (8) shows that C may be put in the form DT with \(D \in D_{+}^{*}\) , \(T \in \mathrm{T}_{1}(n, \mathrm{K})\) , D and T depending continuously on C. The uniqueness of the decomposition X = UDT follows from the uniqueness property of Lemma 10.

The homeomorphism of Prop. 7 is called the Iwasawa decomposition of \(\mathbf{GL}(n,\mathbf{K})\) .

The group \(\mathbf{G} = \mathbf{D}_{+}^{*} \cdot \mathbf{T}_{1}(n, \mathbf{K})\) is the set of upper triangular matrices over K whose diagonal elements are \textgreater0. Let us identity the element \((z_{ij})\) of this group with the element

\[
\left(\left(z _ {i i}\right) _ {1 \leqslant i \leqslant n}, \left(z _ {i j}\right) _ {1 \leqslant i <   j \leqslant n}\right) \in \left(\mathbf {R} _ {+} ^ {*}\right) ^ {n} \times \mathrm{K} ^ {n (n - 1) / 2}.
\]

Arguing exactly as in Example 4, one finds as right Haar measure on this group the measure (when \(\mathbf{K} = \mathbf{R}\) )

\[
\left(\prod_ {i = 1} ^ {n} z _ {i i} ^ {- i}\right) \cdot \left(\bigotimes_ {i = 1} ^ {n} d z _ {i i}\right) \otimes \left(\bigotimes_ {i <   j} d z _ {i j}\right).
\]

Then applying Prop. 13 of §2, No.~9, one sees that if \(\mathbf{GL}(n,\mathbf{K})\) is identified with \(\mathbf{U}(n,\mathbf{K})\times\mathbf{G}\) by the mapping \((U,S)\mapsto US\) , a Haar measure on \(\mathbf{GL}(n,\mathbf{K})\) is given by (when \(\mathbf{K}=\mathbf{R}\) )

\[
\left(\prod_ {i = 1} ^ {n} z _ {i i} ^ {- i}\right) \cdot \alpha \otimes \left(\bigotimes_ {i = 1} ^ {n} d z _ {i i}\right) \otimes \left(\bigotimes_ {i <   n} d z _ {i j}\right),\tag{14}
\]

where \(\alpha\) denotes a Haar measure on \(\mathbf{U}(n,\mathbf{K})\) .

Example 8. --- Spaces of hermitian forms.

In this example, K always denotes one of the fields R, C, H. We write \(\delta = \dim_{R} K\) (thus \(\delta = 1, 2\) or 4). A hermitian form \(\Phi\) on the right vector space \(K^{n}\) may be written

\[
\Phi (x, y) = \Phi (x _ {1}, \dots , x _ {n}, y _ {1}, \dots , y _ {n}) = \sum_ {i, j = 1} ^ {n} \overline {{{{x}}}} _ {i} h _ {i j} y _ {j}
\]

with \(h_{ij} = \overline{h}_{ji}\) for all i and j. We denote by H the vector space over R formed by the hermitian matrices of \(\mathbf{M}_{n}(\mathbf{K})\) . The mapping \((h_{ij}) \mapsto \Phi\) is an isomorphism of H onto the vector space of hermitian forms on \(K^{n}\) , by means of which we shall identify these two spaces. Let \(H_{+}^{*} \subset H\) be the set of nondegenerate positive hermitian forms on \(K^{n}\) . The set \(H_{+}^{*}\) is convex in H; for, if \(\Phi_{1}, \Phi_{2}\) are in \(H_{+}^{*}\) and if \(\lambda, \mu\) are two numbers \textgreater0 such that \(\lambda + \mu = 1\) , it is clear that \(\lambda \Phi_{1} + \mu \Phi_{2}\) is a positive hermitian form; on the other hand, if \((\lambda \Phi_{1} + \mu \Phi_{2})(x, x) = 0\) , then \(\Phi_{1}(x, x) = \Phi_{2}(x, x) = 0\) , therefore x = 0, so that \(\lambda \Phi_{1} + \mu \Phi_{2}\) is nondegenerate. Let us now show that \(H_{+}^{*}\) is an open subset of H. Let S be the set of \(x = (x_{1}, \ldots, x_{n}) \in K^{n}\) such that \(x_{1} \overline{x}_{1} + \cdots + x_{n} \overline{x}_{n} = 1\) ; this is a compact subset of \(K^{n}\) ; if \(\Phi \in H_{+}^{*}\) , the function \(x \mapsto \Phi(x, x)\) is continuous and \textgreater0 on S, hence its infimum is \textgreater0; if \(\Phi' \in H\) is sufficiently near \(\Phi\) , it follows that \(\Phi'(x, x) > 0\) for all \(x \in S\) , so that \(\Phi'\) is positive and nondegenerate.

The general linear group \(\mathbf{GL}(n,\mathbf{K})\) operates continuously on the right in \(\mathfrak{H}\) by \((X,\Phi)\mapsto\Phi\circ X\) , that is, by \((X,H)\mapsto{}^{t}\overline{X}\cdot H\cdot X\) , where \(H\) denotes the hermitian matrix corresponding to \(\Phi\) . It is clear that \(\mathfrak{H}_{+}^{*}\) is stable under \(\mathbf{GL}(n,\mathbf{K})\) . More precisely, by Alg., Ch. IX, §6, No.~1, Cor. 1 of Th. 1, \(^{1}\) \(\mathfrak{H}_{+}^{*}\) is the orbit under \(\mathbf{GL}(n,\mathbf{K})\) of the form \(\sum_{i=1}^{n}\overline{x}_{i}y_{i}\) corresponding to the identity matrix \(I_{n}\) . The stabilizer of this form is \(\mathbf{U}(n,\mathbf{K})\) . By Lemma 2 of App. I, \(\mathfrak{H}_{+}^{*}\) may be identified, as a topological homogeneous space, with \(\mathbf{GL}(n,\mathbf{K})/\mathbf{U}(n,\mathbf{K})\) .

For every \(X \in \mathbf{GL}(n, \mathbf{K})\) , let \(\widetilde{X}\) be the automorphism \(H \mapsto {}^t \overline{X} \cdot H \cdot X\) of the real vector space \(\mathfrak{H}\) . If \(\mu\) denotes the Haar measure of the additive group \(\mathfrak{H}\) , one has \(\widetilde{X}^{-1}(\mu) = |\det \widetilde{X}| \cdot \mu\) (§1, No.~10, Cor. 1 of Prop. 15). Let us show that

\[
| \det \widetilde {X} | = | N (X) | ^ {\lambda},\tag{15}
\]

where N denotes the norm in \(\mathbf{M}_{n}(\mathrm{K})\) regarded as an R-algebra, and where \(\lambda = 1 - \frac{\delta - 2}{\delta n}\) . It suffices to verify (15) for X running over a system of generators of \(\mathbf{GL}(n,\mathbf{K})\) , hence (A, II, §10, No.~13, Cor. 2 of Prop. 14) for \(X\) of the following types:

\begin{enumerate}
\def\labelenumi{\alph{enumi})}
\tightlist
\item
  \(X\) is the matrix of a mapping of the form
\end{enumerate}

\[
\left(x _ {1}, \dots , x _ {n}\right) \mapsto \left(x _ {\sigma (1)}, \dots , x _ {\sigma (n)}\right),
\]

where \(\sigma \in \mathfrak{S}_n\) . In this case, a power of \(X\) is equal to 1, therefore \(|\det \widetilde{X}| = |\mathrm{N}(X)| = 1\) .

\begin{enumerate}
\def\labelenumi{\alph{enumi})}
\setcounter{enumi}{1}
\tightlist
\item
  \(X\) is the matrix of a mapping of the form
\end{enumerate}

\[
\left(x _ {1}, \dots , x _ {n}\right) \mapsto \left(a x _ {1}, x _ {2}, \dots , x _ {n}\right).
\]

Then, if \((h_{ij})\in \mathfrak{H}\) , one has \(\widetilde{X} ((h_{ij})) = (h_{ij}')\) with \(h_{11}' = \overline{a} h_{11}a = |a|^2 h_{11}\) , \(h_{1i}' = \overline{a} h_{1i}\) for \(i > 1\) , \(h_{ij}' = h_{ij}\) for \(i > 1\) , \(j > 1\) ; therefore

\[
| \det \widetilde {X} | = | a | ^ {2} | a | ^ {\delta (n - 1)} = | a | ^ {2 + \delta (n - 1)}.
\]

On the other hand, if \(Y = (y_{ij}) \in \mathbf{M}_n(\mathbf{K})\) , then \(XY = (y_{ij}')\) with \(y_{1j}' = ay_{1j}\) and \(y_{ij}' = y_{ij}\) for \(i > 1\) ; therefore \(|\mathrm{N}(X)| = |a|^{\delta n}\) . The formula (15) is again verified.

\begin{enumerate}
\def\labelenumi{\alph{enumi})}
\setcounter{enumi}{2}
\tightlist
\item
  \(X\) is the matrix of a mapping of the form
\end{enumerate}

\[
\left(x _ {1}, \dots , x _ {n}\right) \mapsto \left(x _ {1} + b x _ {2}, x _ {2}, \dots , x _ {n}\right).
\]

Then \(\widetilde{X}\big((h_{ij})\big) = (h_{ij}^{\prime})\) with \(h_{11}^{\prime} = h_{11}\) , \(h_{12}^{\prime} = h_{12} + h_{11}b\) , \(h_{1i}^{\prime} = h_{1i}\) for \(i > 2\) , \(h_{22}^{\prime} = h_{22} + \overline{b}h_{12} + \overline{h}_{12}b + \overline{b}h_{11}b\) , \(h_{2i}^{\prime} = h_{2i} + \overline{b}h_{1i}\) for \(i > 2\) , \(h_{ij}^{\prime} = h_{ij}\) for \(i > 2\) , \(j > 2\) . Taking into account Lemma 6, one sees that \(|\det \widetilde{X}| = 1\) . One verifies similarly that \(|\mathrm{N}(X)| = 1\) , which completes the proof of formula (15).

This established, the measure \(|\mathrm{N}(H)|^{-\lambda/2}d\mu(H)\) on H is invariant under \(\mathbf{GL}(n,\mathbf{K})\) , since

\[
\begin{array}{r l} & {\widetilde {X} ^ {- 1} \big (| \mathrm{N} (H) | ^ {- \lambda / 2} d \mu (H) \big) = \big | \mathrm{N} \big (\widetilde {X} (H) \big) \big | ^ {- \lambda / 2} \cdot | \det \widetilde {X} | d \mu (H)} \\ & {\qquad = | \mathrm{N} (H) | ^ {- \lambda / 2} | \mathrm{N} (X) | ^ {- \lambda} | \det \widetilde {X} | d \mu (H) = | \mathrm{N} (H) | ^ {- \lambda / 2} d \mu (H).} \end{array}
\]

If \(H \in \mathfrak{H}_+^*\) , then \(H = \widetilde{X}(\mathrm{I}_n) = {}^t\overline{X}X\) for some \(X \in \mathbf{GL}(n,\mathbf{K})\) , therefore \(\mathrm{N}(H) = \overline{\mathrm{N}(X)}\mathrm{N}(X) > 0\) . Consequently, on \(\mathfrak{H}_+^*\) , the unique (up to a constant factor) measure invariant under \(\mathbf{GL}(n,\mathbf{K})\) (cf.~§2, No.~6, Th. 3) is the measure

\[
d \gamma (H) = \mathrm{N} (H) ^ {- \lambda / 2} d \mu (H).
\]

In particular,

\[
d \gamma (H) = (\det H) ^ {- (n + 1) / 2} d \mu (H) \quad \text { when } K = R,
\]

\[
d \gamma (H) = (\det H) ^ {- n} d \mu (H) \quad \text { when } \mathrm{K} = \mathrm{C}.
\]

